\chapter{Introduction}

% Accroche & introduction du sujet

\begin{quote}
    \textit{Mais dès l'instant où le critère de l'authenticité s'avère inapplicable à la production artistique, la fonction sociale entière de l'art s'en trouve bouleversée. Au lieu de se fonder sur le rituel, elle se fonde sur une autre praxis : en l'occurrence, sur la politique.}
    
    Walter Benjamin, \textit{L'\oe uvre d'art à l'époque de sa reproductibilité technique}, 1939.
\end{quote}

L'acte de copie d'un objet culturel est la première étape à son partage, sa conservation et sa réinterprétation. Si une langue a besoin de locuteurs pour pouvoir évoluer, alors un objet culturel qui n'est pas copié ne peut pas être nommé comme un \guillemetleft objet culturel \guillemetright ~, tout comme le latin ne peut pas être considéré comme une \guillemetleft langue vivante \guillemetright. Un premier exemple de l'importance de la copie et par extension du piratage dans notre patrimoine est \textit{Le Voyage dans la Lune} réalisé en 1902 par Georges Méliès. La plupart des copies officielles de ce film ont disparu, ou ont été détruites par Méliès lui-même alors qu'il ne pouvait plus les entreposer à cause de la fermeture de son studio. Les chercheurs et restaurateurs ont pu reconstituer et préserver son œuvre grâce à de très nombreuses copies qui étaient soit piratées, soit confisquées \footnote{Marie-France Briselance et Jean-Claude Morin, \textit{Grammaire du cinéma}, Paris, Nouveau Monde, coll. « Cinéma », 2010, 588 p. (ISBN 978-2-84736-458-3), p. 111.}.

Un deuxième exemple, peut-être le plus connu, est celui du film \textit{Nosferatu le vampire} réalisé en 1922 par Friedrich Wilhelm Murnau. Florence Stoker, veuve de Bram Stoker, l'auteur du roman \textit{Dracula}, intente à Murnau un procès pour plagiat que ce dernier perd. Toutes les copies ainsi que les négatifs sont détruits par deux fois en 1925 et 1928. Rapidement, d'anciennes copies cachées refont surface - et le film survit grâce à des projections confidentielles jusqu'en 1981 où plusieurs copies partielles permettent de le restaurer au plus près de l'archétype du travail de Murnau \footnote{Jean Pierre Piton, \textit{L'éternel retour de Dracula}, L'Écran Fantastique no 130 janvier 1993 p. 21.}.

La numérisation de nos outils professionnels, de nos sources de divertissement et de nos modes de communication, généralise la capacité de copie des objets culturels qui peuplent nos vies. Cette propriété préexistait déjà avec, par exemple, la K7 VHS à bande magnétique qui permet d'enregistrer un signal analogique. Cependant, dans un objet numérique, la propriété de copie est parfaitement intégrée au fonctionnement du fichier et de son système d'accès par défaut, l'ordinateur, produisant un dilemme numérique entre la liberté du partage d'une information réplicable gratuitement et la propriété intellectuelle littéraire et artistique \footnote{\cite{national_research_council_us_digital_2010}}.

La nouvelle facilité de copie et d'échange de l'information grâce à la généralisation de l'informatique bon marché et de l'accès au web domestique a rapproché des individus qui se sont organisés en communautés culturelles,  partageant des références culturelles et les fichiers qui les documentent, devenant ainsi des communautés patrimoniales.

Les communautés patrimoniales \footnote{~Convention de Faro sur la valeur du patrimoine culturel pour la société, 2005} qui nous intéressent dans ce travail sont celles qui décident d'embrasser la copie et le partage de l'information dans sa forme la plus radicale, en privilégiant les \guillemetleft communs numériques \guillemetright. Elles peuvent choisir de produire et d'entretenir leur propre contenu à la manière d'un Wikipédia, devenant des \textit{communautés libres}, ou elles peuvent choisir de faire fi des règles et des lois en devenant des \textit{communautés patrimoniales pirates}\footnote{\cite{bermes_lecran_2024}}. Si les communautés libres essaient de se différencier le plus possible des communautés pirates, les individus les composant évoluent entre les deux, puisant dans les ressources des communautés libres pour ensuite compléter leurs besoins avec l'aide du piratage.

Nous étudierons ici les liens, les émotions, et les mouvements de va-et-vient entre les individus qui composent les communautés pirates, les objets numériques piratés, et les lieux virtuels où se construisent ces communautés. Notre premier but est d'étendre le concept de communautés patrimoniales aux pirates culturels, puis de comprendre l’appréhension du public du \guillemetleft dilemme numérique de la propriété intellectuelle \guillemetright, et enfin d'étudier les pratiques émergentes de partage et de conservation numériques qui mobilisent ces communautés patrimoniales pirates.

\section{Langues et jambes de bois : définir la piraterie culturelle}

Les travaux sur la véritable piraterie historique ont souvent une introduction permettant de séparer une partie des représentations du sujet et la réalité des faits. Je pense que pour la piraterie culturelle, un travail similaire est le bienvenu. Les travaux sur le piratage se répartissent sur deux axes bibliographiques principaux.\\

\subsection*{La piraterie culturel entre les enjeux scientifiques et la communication}

Les travaux sur le piratage se répartissent sur deux axes bibliographiques principaux.\\

Le premier se concentre sur les industries de la propriété intellectuelle, essentiellement dans le domaine du marché pharmaceutique international. Pour ce qui est de la piraterie de biens culturels, le sujet le plus étudié est la place du piratage dans les économies des pays en développement, à la fois ateliers du monde du capitalisme planétaire, mais aussi envisagé comme une manne de consommateurs potentiels des divertissements du nord économique. Des revues de science de l'information, de communication, de management et d'économie étudient l'organisation et les chemins de la distribution du contenu culturel des pays du Sud économique comme par exemple à Malégaon en Inde \footnote{Lawrence Liang,\textit{ Piratage, créativité et infrastructure : repenser l’accès à la culture}, Tracés. Revue de Sciences humaines, 2014}, ou au Cameroun \footnote{Mattelart, T.\textit{ Chapitre 1. Piratages audiovisuels et réseaux de la mondialisation par le bas}. Dans \og Piratages audiovisuels : Les voies souterraines de la mondialisation culturelle \fg, p. 27-52, éditions De Boeck Supérieur, 2011}.

Le deuxième se concentre sur la piraterie culturelle numérique comme un marché, où le monde académique tente de saisir à la fois les lieux où s'organisent les échanges, et de comprendre comment s'organisent ceux qui font vivre ces lieux. Le mot de \guillemetleft  Pirate \guillemetright~, ainsi que l'iconographie de la piraterie historique, est également un point de rendez-vous, ou un repoussoir - pour les communautés qui investissent le web comme un espace d'échange de biens culturels. L'exemple le plus partagé est le nom et l'iconographie du site \textit{The Pirate Bay}, dont l'icône est un trois-mâts noir, dont la voile centrale montre une variation du \textit{Jolly Roger} remplaçant le crâne par une cassette en inscrivant en dessous dans une police gothique le nom du site de partage de torrents le plus connu des internets. L'utilisation d'autres termes comme \textit{partage illégal du contenu culturel, reproduction interdite ou non autorisée, partage d'informations avec son voisin} est une alternative qui offre une représentation plus mesurée et plus proche de la réalité des faits \footnote{Richard Stalman,\textit{Quelques expressions trompeuses qu'il vaut mieux éviter d'employer}, dans \guillemetleft Libre enfants du savoir numérique \guillemetright, p.359-364 }. Au-delà de cette analyse légaliste et quantitative, la sphère académique des sciences humaines et sociales s'est intéressée à la popularité de plateformes telles que Napster \footnote{Moulier-Boutang, \textit{Droits de propriété intellectuelle, terra nullius et capitalisme cognitif}. Multitudes,n°41(2), p66-72, 2010} ou The Pirate Bay . Ces recherches universitaires ont pour but de comprendre la psychologie des individus en examinant les discours et valeurs associés au piratage.

\subsection*{La communauté patrimoniale pirate}

L'objet numérique et son partage mobilisent des concepts et des catégories juridiques diverses qu'il est important de définir. Cet objet numérique s'incorpore dans le concept de \textit{patrimoine documentaire immatériel}, promu et adopté par la Convention de l'UNESCO pour que ses États membres puissent intégrer au droit des définitions et des outils politiques de promotion de la diversité culturelle, y compris de la culture numérique et de la culture du web. Si un patrimoine est immatériel, c'est alors qu'il est constitué d'objets qui ont un sens particulier pour un groupe de personnes, le patrimoine immatériel étant donc la relation entre ces \textit{objets} et les groupes de personnes qui leur donnent un sens et une fonction. Quand un objet est la cause principale de lien dans un groupe (et non plus une conséquence de ce lien), nous pouvons alors parler de l'apparition d'une \textit{communauté patrimoniale}. Ce terme de \textit{communauté patrimoniale} entre dans le vocabulaire institutionnel en 2005 avec la signature de la convention européenne de Faro sur la valeur du patrimoine culturel pour la société. Cette convention inscrit l'action européenne dans le domaine du patrimoine avec un nouveau paradigme, considérant autant les objets et les lieux que les usages, les personnes et les interactions qui en résultent.

Les émotions négatives sont une composante puissante dans la définition de ce qu'est le patrimoine (peur de voir disparaître, de voir changer, de voir trahir un objet), les communautés se soudent avec des émotions comme la colère, la sédition, la rébellion face à l'injustice \footnote{\cite{fabre_emotions_2013}}, en invoquant des référentiels politiques ou institutionnels.  Les communautés patrimoniales peuvent entrer dans des dynamiques d'actions de sauvegarde de leur patrimoine, parfois avec la volonté d'aider les acteurs institutionnels ou entrepreneuriaux (une restauration d'un bâtiment par exemple), ou en entrant dans un rapport de force avec ces mêmes acteurs (la mobilisation des mouvements écologistes face à des travaux d'infrastructures comme les bassines de Sainte-Soline en est un bon exemple). 

Si le partage non autorisé de biens culturels n'a rien de nouveau, le tournant numérique de nos sociétés a engendré deux conséquences. La première est la massification des échanges, des capacités de stockage et de copie des individus,  la deuxième est la capacité nouvelle des individus à s'affranchir du temps et de l'espace pour entrer en relation, donnant ainsi naissance à des communautés patrimoniales plus nombreuses et beaucoup plus développées dès la démocratisation de l'accès au web. Quand une communauté patrimoniale se forme autour d'objets ou de pratiques \textit{illégales} nous pouvons alors utiliser le terme de communautés patrimoniales \textit{pirates}. Notre travail a pour objectif de mieux comprendre les relations entre l'objet, l'individu et la communauté dans ce cadre de sortie des référentiels institutionnels et entrepreneuriaux - le but des pirates n'étant pas de trouver une légitimité aux yeux de l'État ou des entreprises, mais bien de s'en affranchir.

\section{Revue bibliographique}

Étudier des individus qui piratent implique de se plonger dans leurs histoires personnelles, et les origines de leur montée en compétence progressive - passant de néophyte à amateur, puis pour certains à des professionnels du piratage, ce qu'Anthony Galluzzo, professeur en science de gestion à l'université de Saint-Étienne, nomme une \textit{carrière pirate}\footnote{\cite{galluzzo_longevite_2021}}. Cette acquisition de compétences techniques est parallèle à une éducation, un apprentissage social et politique. C'est en combinant les deux que les pirates génèrent une sociabilité particulière s'appuyant autant sur des éléments techniques (comment accéder à une ressource, qui possède un tutoriel pour tel outil logiciel, qui possède une copie d'un objet numérique rare que je cherche) que sur des émotions patrimoniales (ce site va disparaître ; comment aider ma communauté à le préserver). 

Ce lien a été l'objet d'étude de la littérature anglo-saxonne des \textit{fan studies} \footnote{\cite{taylor_save_nodate}} \footnote{:\cite{de_kosnik_rogue_2016}}. Ces études se concentrent sur les cultures du web (des fanfictions, le jeu vidéo) et la manière dont les personnes qui vivent ces cultures comprennent et organisent la mémoire de ce qui a été produit au sein de leur communauté. Nous étudions une culture \textit{nativement numérique}, mais certaines différences sur la situation des communautés patrimoniales pirates ne permettent pas de plaquer les analyses des \textit{fantstudies} sur ce sujet.

L'un des paramètres ayant le plus d'influence sur l'organisation des communautés patrimoniales pirates et de leurs membres est \textit{l'organisation anarchique de guérilla}. Comme pour Wikipédia, la plupart des dépôts communautaires pirates n'obéissent pas à une structure verticale et rigide. Un noyau dur de membres reconnus pour leur implication et leur ancienneté dans le projet permet de maintenir l'infrastructure de base du dépôt communautaire, ce à quoi s'ajoutent des membres moins investis qui offrent leur aide et leurs conseils de manière ponctuelle. Un troisième cercle est la communauté large autour du dépôt, des usagers réguliers offrant un soutien financier, qui propagent la bonne adresse à l'aide du bouche-à-oreille sur les forums, ou simplement qui témoignent de l'importance du dépôt pirate dans leur vie de tous les jours. 

Le terme de guérilla fait référence à la manière dont ces dépôts, leurs membres et leurs usagers s'organisent pour mettre en place un espace à la fois efficace dans sa mission (transmettre et archiver des données) sans se faire repérer par les éditeurs ou les autorités de protection des droits d'auteur \footnote{Martin Paule Eve, Lessons from the Library: Extreme Minimalist Scaling at Pirate Ebook Platforms, dans \textit{Digital Humanities Quarterly} Vol 16 n°2, 2022}. Pour cela, les dépôts pirates doivent à la fois avoir une visibilité minimale (ne pas être vus), sans que le trafic internet puisse être retracé à l'hébergeur du dépôt (ne pas être visé), avec une résilience aux attaques et aux fermetures d'hébergeur (ne pas être touchés).

\section{Nos sources}

Une partie de notre étude se base sur l'étude des archives et des forums de communautés patrimoniales pirates. Une partie a fait l'objet d'un archivage, parfois d'une patrimonialisation et d'une curation, c'est par exemple le cas des archives du web de la BnF ou des collections de The Internet Archive. Une autre partie de nos sources a été directement archivée en utilisant l'outil d'instantané (snapshot) Zotero.
\begin{itemize}
    \item Nous avons effectué un travail de recherche sur les archives du web de la BnF, en particulier les archives du site Skyblog. Grâce à Marina Hervieu, chargé de projet de recherche à la BnF, nous avons pu explorer les skyblogs ayant dans leurs pseudonymes des mots laissant penser à une pratique régulière du piratage \footnote{pirate, piratage, téléchargement, télécharger, torrenting, torrent, fansub, mininova, des mots ou des noms de service en lien avec le piraterie durant les années de service de Skyblog}. Dans les listes de blog minées par Mme Hervieu, nous avons ensuite pu appliquer un tri en fonction de la popularité et du nombre de posts du blog. En utilisant cette méthode, nous avons très rapidement atteint une saturation dans nos d'exemples. 
    
    \item Une grande proportion de nos sources provient de l'observation des archives et de l'actualité des forums des communautés patrimoniales pirates, notamment sur Reddit, Twitter/X et Discord. Les deux sites les plus propices à l'étude se sont révélés être Discord et Reddit, du fait que la plateforme Twitter/X n'a introduit la possibilité de former des « communautés » que très récemment. 
    
    \subitem Discord est un système de messagerie communautaire où nous avons suivi la communauté qui se nomme \textit{The Eye}. C'est un site communautaire issu de la communauté \textit{Archive Team}, ayant pour but de préserver et de partager des pans de la culture numérique. 

    \begin{quote}
        \textit{Living without the knowledge of our past history, origin and culture is like a tree without roots. The Internet is a worldwide platform for sharing information. It is a community of common interests.
        Our mission at The Eye is to preserve pieces of digital history. We are digital librarians.}
        
        Message de présentation du serveur discord de \textit{The Eye}
    \end{quote}

    \subitem le site Reddit est également une source très instructive pour étudier les communautés patrimoniales pirates. Il est composé de \textit{subreddits}, des forums communautaires sur des sujets spécifiques. Il existe plusieurs subreddit sur les sujets du piratage de bien culturel et de la sauvegarde du patrimoine numérique. Bien que l'intégralité d'un subreddit ne puisse être classifiée comme une communauté patrimoniale, certaines communautés investissent activement des subreddits spécifiques. Nous pouvons illustrer cela par \textit{r/datahoarder} (des individus collectant de vastes quantités de données numériques) qui sert de point central pour les discussions et la mobilisation des amateurs d'archives numériques et de patrimoine dématérialisé, ou r/Piracy qui sert de point de ralliement pour toutes les discussions concernant le piratage.

    \subitem Twitter/X est une source plus difficilement exploitable pour notre travail car il s'agit d'une application de microblog où les gens peuvent se suivre de manière individuelle sans avoir l'option d’agréger les utilisateurs en communautés\footnote{Il existe depuis la reprise de l'application par Elon Musk un onglet communauté, mais la perte de popularité de Twitter couplé à l'ajout de cette fonction relativement peu utilisée fait qu'il reste une application de microblog individualisée. }. Les sources que nous agrégeons de Twitter ne sont pas représentatives d'une communauté patrimoniale, cependant certains sujets transversaux peuvent mobiliser un grand groupe d'utilisateurs, créant ainsi un effet de mode et d’entraînement permettant d'étudier l'importance de certains évènements sur la création de communautés patrimoniales pirates. L'exemple le plus récent dont je peux témoigner est la mort du réalisateur américain David Lynch. Les \textit{twittos}\footnote{Twittos : surnom vulgaire donné à un utilisateur intensif du réseau Twitter/X} et fans endeuillé.es ont commencé à partager illégalement les films et séries et documentaires du réalisateur via des services comme Mega ou Google Drive. 
    
    \end{itemize}

Dans le temps qui nous a été imparti, nous avons également réalisé cinq entretiens. Notre but était de pouvoir fournir une diversité d'âge, de genre, de pratiques et d'expériences. Ces entretiens ont été pour la plupart réalisés en tête-à-tête, sauf pour Théo qui habite trop loin. Une transcription automatique a été systématiquement mise en place par captation audio et traitement utilisant l'outil Whisper d'OpenAI.\newline

Tableau des profils des différents entretiens.\\
\begin{tabular}{|p{2cm}|p{2cm}|p{2cm}|p{2cm}|p{2cm}|}

\hline
   Prénom  & Âge & Profession & Méthode & Résumé\\
\hline
    Adrien & 23 ans & Étudiant en informatique & Torrenting, Téléchargement & Diplôme en informatique \\
\hline
    Théo & 27 ans & Ingénieur informaticien & Torrenting, Serveur & Diplôme en informatique\\
\hline
    Hilaire & 26 ans & Étudiant en sociologie & Téléchargement, forums, Torrenting & Diplôme en communication et en sociologie \\
\hline 
    Zakary & 25 ans & Étudiant en communication & Téléchargement, communauté & Diplôme en histoire et communication \\
\hline
    Elise & 23 ans & Étudiante en communication & Téléchargement, communauté & licence langue et édition, master de communication \\
\hline
\end{tabular}\\

Nous avons tenté de représenter équitablement les genres, mais un temps de recherche insuffisant entraîne un biais persistant lié à l'âge. La plupart des répondants sont des personnes de ma génération (nées entre 1995 et 2003). Nous avons cependant réussi à mêler des profils de personnes proches de l'informatique (la majorité du corps social qui m'entoure) et d'autres qui piratent sans pouvoir se reposer sur des connaissances techniques aussi poussées. 
Malgré une cohorte limitée, la saturation sémantique des réponses de nos enquêtés a été relativement rapide. La plupart de nos répondants ont une pratique du piratage privilégiant l'accès le plus simple et le plus rapide aux ressources, homogénéisant rapidement les réponses aux questions. Les entretiens individuels sont difficilement des portes d'entrée vers les communautés patrimoniales pirates - qui sont par définition des œuvres collectives ; ils restent cependant très intéressants pour comprendre la manière dont des individus interagissent avec ces communautés sans pour autant en faire partie.

Nous avons également utilisé comme sources plusieurs bibliothèques pirates : Sci-Hub, LibGen et Zlibrary ; pour comprendre et étudier à la fois la culture pirate, les outils et l'expérience que les individus trouvent sur ces bibliothèques, ainsi que l'architecture et l'organisation de ces sites. 

\section{Méthodologie et problématique}

%méthodologie

Notre méthodologie s'inspire à la fois des \textit{fanstudies}, de la sociologie de la déviance et des études en archivistique et en sciences de l'information et des bibliothèques. Notre but est de rapprocher les expériences personnelles des personnes qui piratent et l'étude de l'architecture organisationnelle des bibliothèques pirates pour comprendre les éléments et les formes d'organisation dans les communautés patrimoniales pirates. Nous essaierons de comprendre et de cartographier les émotions patrimoniales appliquées à des objets culturels numériques, et la manière dont ces émotions sont le fondement de l'émergence d'une pratique archivistique populaire. Notre sujet d'observation sera les dépôts pirates, des bibliothèques et archives communautaires, autogestionnaires et clandestines. Nous étudierons à la fois l'aspect sensible et intime de l'émotion patrimoniale, du vécu pirate, mais également le contexte d'organisation et de remise en question des normes de distribution et de conservation de la culture sur Internet. Nous cherchons à explorer la culture et les références qui permettent à des individus qui ne se sont sans doute jamais vus de s'organiser, et d'impacter le circuit de la diffusion de l'information en cherchant à inventer des modèles alternatifs de gestion et de partage du patrimoine immatériel. \\

%annonce de plan 

D'abord, nous nous concentrerons sur la vie intérieure des pirates, de leurs expériences et l'apprentissage social qui les sépare des communautés patrimoniales pirates. Ensuite, nous aborderons la culture pirate, sa vision du monde et son lien avec l'émotion patrimoniale de sédition et les problématiques d'organisation qui peuvent en ressortir. Enfin, nous irons nous immerger dans la grande bibliothèque, en édifiant une typologie des différents types de dépôts pirates, de leurs organisations et des différents buts pour lesquels ils ont été créés.

%Nous nous essaierons dans ce travail de comprendre les dynamiques culturelles et les organisations sociales sous-jacentes aux pratiques des communautés patrimoniales pirates. 
%Comment les dynamiques culturelles et les organisations sociales permettent-elles l'émergence d'une pratique archivistique populaire, autogestionnaire et clandestine ? Notre recherche s'inscrit dans un contexte de remise en question international des normes de distribution et de conservation du savoir, posant de nouveaux défis aux institutions patrimoniales et aux acteurs entrepreneuriaux de la distribution culturelle. Nous cherchons à explorer la manière dont ces entités informelles s'organisent, pour comprendre l'impact sur les circuits de la diffusion de la culture, mais également pour envisager des modèles alternatifs de gestion et de partage du patrimoine immatériel.\\

%Nous étudierons d'abord le chemin individuel d'apprentissage, de partage et de mise en relation des individus avec les communautés patrimoniales pirates. Ensuite, nous aborderons la vision du monde des individus qui piratent et les liens entre émotions pirates et émotions patrimoniales. Enfin, nous irons visiter plusieurs dépôts numériques pirates marqués par les personnalités et intérêts des individus qui les composent. 


